\section{Introduction} \label{sec1}
It was noted by Fock \cite{Fock:1937} in 1937 that it is convenient
to represent Green functions as integrals over an auxiliary
coordinate (the so-called ``proper time'') of a kernel satisfying
the heat equation. Later on Schwinger \cite{Schwinger:1951} recognised 
that this representations makes many issues related to renormalization
and gauge invariance in external fields
more transparent. These two works introduced the
heat kernel to quantum theory. DeWitt made the heat kernel one of the
main tools of his manifestly covariant approach 
\cite{DeWitt:1965,Dewitt:1967yk,Dewitt:1967ub,Dewitt:1967uc}
to quantum field theory and quantum gravity which became dominant 
for many years. 

Heat kernel is a classical subject in mathematics 
\cite{Hadamard:1932}\footnote{A historical survey of 
the mathematical literature on the heat kernel expansion
can be found in \cite{Gilkey:1995}.}.
Asymptotics of the heat kernel are closely related to the eigenvalue
asymptotics  found by H.~Weyl \cite{Weyl:1912,Weyl:1915} 
and studied further in \cite{Mina,MinaPle}. The problem, as it was
formulated by Kac \cite{Kac:1966}, reads: ``Can one hear the shape of
a drum?''. In other words, this is the problem of recovering geometry
of a manifold from the spectrum of a natural differential operator.
Heat kernel coefficients proved very useful in this context.
On the other hand, the heat kernel is also an adequate tool
to study the index theorem of Atiyah and Singer 
\cite{Atiyah:1963a,Gilkey:1973,AtiyahBP:1973}.

By about 1990 the heat kernel expansion on manifolds without boundaries
or with boundaries and simplest local boundary conditions on them
was well understood. Also, the heat kernel became a standard tool
in calculations of the vacuum polarisation, the Casimir effect, and
in study of quantum anomalies. Later on, progress in theoretical
physics, especially in string theory and related areas, and parallel
developments in mathematics made this field highly specialised.
New results on non-standard boundary conditions (as, e.g., containing
tangential derivatives on the boundary or non-localities) or on
non-standard geometries (domain walls) were scattered in large
amounts of physical and mathematical literature. The aim of this
report is to present a unifying approach to the heat kernel expansion
and to supply the reader with a ``user-friendly'' guide to
the field. The main idea which we shall pursue is the universality of
the heat kernel. A single calculation (though, sometimes, quite
involved) may help in a large variety of applications regardless
of such details as spin, gauge group, etc. As well, just a single
universal object in fact describes counterterms, anomalies,
some asymptotics of the effective action, and much more.

To illustrate the use of the heat kernel in quantum field theory
let us
consider the generating functional for the Green functions of the 
field $\phi$ in the path integral representation
\begin{equation}
Z[J]=\int \mathcal{D} \phi \exp (-\mathcal{L}(\phi,J)) \ . \label{pathint}
\end{equation}
The heat kernel methods are almost exclusively
used for the one-loop calculations. In this approximation it is enough
to expand the action $\mathcal{L}$ to up to the quadratic order in quantum 
fluctuations $\phi$.
\begin{equation}
\mathcal{L}=\mathcal{L}_{\rm cl} + \langle \phi ,J \rangle +
\langle \phi, D\phi \rangle \,,\label{actquad}
\end{equation}
where $\mathcal{L}_{\rm cl}$ is the action on a classical background,
$\langle .,.\rangle$ denotes an inner product on the space of quantum
fields. Usually, this inner product is just an integral over the underlying
space. For real one-component scalar fields it reads
\begin{equation}
\langle \phi_1, \phi_2 \rangle =\int d^n x \sqrt{g} \phi_1(x)\phi_2(x)
\,.\label{inprosca}
\end{equation}
The linear term in (\ref{actquad}) contains in general contributions
from the external sources of the field $\phi$ and from the first 
variation of the classical action. If the background satisfies 
classical equation of motion for the field $\phi$, the latter part
of the linear term vanishes, though the former one (external sources)
should be kept arbitrary if one wishes to study correlation functions
of $\phi$. We stress that the background and quantum fields may be of
completely different nature. For example, it is a meaningful problem
to consider pure quantum scalar fields on the background of pure
classical gravity. $D$ is a differential operator. After a suitable 
number of integrations by part it is always possible to convert the
quadratic part of the action to the form given in (\ref{actquad}).
We postpone discussion of possible boundary terms to the next
sections. In the simplest case of quantum scalar field on the background
of a classical geometry
$D$ is a Laplacian with a mass term:
\begin{equation}
D=D_0:=-\nabla_\mu\nabla^\mu +m^2 \,.\label{simplestD}
\end{equation}
Just this simple example is enough to illustrate the material of this
section. Note, that in this case $J$ has no contribution from
the first variation of the classical action since gravity is not
quantised.


The path integral measure is defined by:
\begin{equation}
1=\int \mathcal{D} \phi \exp (-\langle\phi ,\phi \rangle) \ . \label{measure}
\end{equation}
Strictly speaking the right hand side of (\ref{measure}) is divergent.
The essence of the condition (\ref{measure}) is that
this divergence does not depend on external sources and on the
background geometry and, therefore, may be absorbed in an irrelevant
normalisation constant. 
The Gaussian integral (\ref{pathint}) can be evaluated giving
\begin{equation}
Z[J]=e^{-\mathcal{L}_{\rm cl}}
{\det }^{-\frac 12} (D) \exp \left( \frac 14 J D^{-1} J \right)
 \, . \label{pathdet}
\end{equation}

We stress that the equation (\ref{pathdet}) is true only if
the operator $D$ is self-adjoint. This means that $D$ is 
symmetric or formally self-adjoint with respect to the
scalar product $\langle .,.\rangle$:
\begin{equation}
\langle \phi_1,D\phi_2 \rangle =
\langle D\phi_1,\phi_2 \rangle \label{def-sym-op}
\end{equation}
for any $\phi_1$, $\phi_2$, and that the domains of definition
of $D$ and its adjoint coincide. We will not care about the
second requirement since it involves mathematical machinery
\cite{ReedSimon}
going beyond the scope of the present report. The first requirement
(\ref{def-sym-op}) poses important restrictions on admissible
boundary conditions.

To become convinced that (\ref{def-sym-op}) is really necessary,
one can calculate a finite-dimensional Gaussian integral
with 
\begin{equation}
D=\left( \begin{array}{cc} a & b \\ c& d \end{array} \right)
\end{equation}
and $a$, $b$, $c$ and $d$ being real constants first by completing the
squares in the exponent, and then compare the result with
$\det D$ (of course, one should remember a factor
of $\pi$). The two results coincide if $b=c$.

Let us return to the generating functional (\ref{pathdet}).
To analyse the two multiplies on the right hand side of
(\ref{pathdet}) which depend on the operator $D$ it is
convenient to introduce the heat kernel
\begin{equation}
K(t;x,y;D)=\langle x| \exp (-tD) | y\rangle \,.
\label{hkmatrel}
\end{equation}
This somewhat formal expression means that $K(t;x,y;D)$ should
satisfy the heat conduction equation 
\begin{equation}
(\partial_t+D_x) K(t;x,y;D)=0 \label{heateq}
\end{equation}
with the initial condition
\begin{equation}
K(0;x,y;D)=\delta (x,y)\,.\label{inicon}
\end{equation}

For $D=D_0$ (\ref{simplestD}) on a flat manifold $M=\mathbb{R}^n$ 
the heat kernel reads:
\begin{equation}
K(t;x,y;D_0)= (4\pi t)^{-n/2} \exp\left(
-\frac{ (x-y)^2}{4t} -tm^2 \right) \,.\label{simplestHK}
\end{equation}
The equations (\ref{heateq}) and (\ref{inicon}) can be checked
straightforwardly. Let us consider a more general operator $D$
which contains also a potential term or a gauge field. Then
(\ref{simplestHK}) still describes the leading singularity
 in the heat kernel as $t\to 0$. The subleading terms have a
form of the power-law corrections\footnote{
On a curved space even the leading term must be modified
(cf sec.\ \ref{s4iter}). On manifolds with boundaries also half-integer
powers of $t$ appear in the expansion, and, consequently,
$b_{2j+1}\ne 0$.}:
\begin{equation}
K(t;x,y;D)=K(t;x,y;D_0) \left( 1 +tb_2(x,y)+t^2b_4(x,y)+\dots
\right) \,.\label{sHKexp}
\end{equation}
The coefficients $b_k(x,y)$ are regular in the limit $y\to x$.
They are called the heat kernel coefficients\footnote{We shall
mostly use the coefficients $a_k$ which differ from $b_k$ by a
normalisation factor.}. At coinciding arguments $b_k(x,x)$ 
are local polynomials of background fields and their derivatives.

The propagator $D^{-1}(x,y)$ can be defined through the heat kernel
by the integral representation
\begin{equation}
D^{-1}(x,y)=\int_0^\infty dt\, K(t;x,y;D)\,,
\label{proptoHK}
\end{equation}
which follows from (\ref{hkmatrel}) if we suppose that the heat kernel
vanishes sufficiently fast as $t\to \infty$. We can formally integrate
the expansion (\ref{sHKexp}) to obtain
\begin{equation}
D^{-1}(x,y) \simeq 2(4\pi )^{-n/2} \sum_{j=0}
\left( \frac{ |x-y|}{2m} \right)^{-\frac 12 n+j+1}
K_{-\frac 12 n+j+1} (|x-y|m)b_{2j}(x,y) \,,\label{DBessel}
\end{equation} 
where $b_0=1$. By examining the behaviour of the Bessel function
$K_\nu (z)$ for small argument $z$ \cite{Gradshteyn:2000}
we conclude that the singularities in the propagator at
coinciding points are described by the first several heat kernel
coefficients $b_k$.

Let us consider the part of the generating functional (\ref{pathdet})
which contains $\det (D)$. The functional 
\begin{equation}
W=\frac 12 \ln \det (D) \label{defW}
\end{equation}
is called the one-loop effective action. It describes the quantum 
effects due to the background fields in the one-loop approximation
of quantum fields theory. To relate $W$ to the heat kernel we shall
use the arguments of \cite{Birrell:1982ix}. 

For each positive eigenvalue $\lambda$ of the operator $D$ we may write
an identity
\begin{equation}
\ln \lambda =-\int_0^\infty \frac {dt}t e^{-t\lambda}\,.
\label{lnlam}
\end{equation}
This identity is ``correct'' up to an infinite constant, which does
not depend on $\lambda$ and, therefore, may be ignored in what 
follows\footnote{To ``prove'' this statement one has to differentiate
both sides of eq.\ (\ref{lnlam}) with respect to $\lambda$.
}. Now we use $\ln\det (D)={\rm Tr}\ln (D)$ and extend
(\ref{lnlam}) to the whole operator $D$ to obtain 
\begin{equation}
W=-\frac 12 \int_0^\infty \frac {dt}t K(t,D) \,,\label{WKtD}
\end{equation}
where
\begin{equation}
K(t,D)={\rm Tr}\left( e^{-tD} \right)=\int d^nx \sqrt{g} K(t;x,x;D)
\,.\label{KtDdef}
\end{equation}

Here we have only presented some heuristic arguments in favour of
eq.\ (\ref{WKtD}). A more rigorous treatment of functional determinants
can be found in sec.\ \ref{s2spec}.

The integral in (\ref{WKtD}) may be divergent at both limits.
Divergences at $t=\infty$ are caused by zero or negative eigenvalues
of $D$. These are the infra red divergences. They will not be discussed
in this section. We simply suppose that the mass $m$ is sufficiently
large to make the integral (\ref{WKtD}) convergent at the upper limit.
Divergences at the lower limit cannot be removed in such a way.
Let us introduce a cut off at $t=\Lambda^{-2}$.
\begin{equation}
W_\Lambda = -\frac 12 \int_{\Lambda^{-2}}^\infty \frac {dt}t K(t,D)
\,.\label{WLamreg}
\end{equation}
It is now easy to calculate the part of $W_\Lambda$ which
diverges in the limit $\Lambda\to\infty$:
\begin{eqnarray}
&&W_\Lambda^{\rm div} =-(4\pi )^{-n/2} \int d^n x \sqrt{g} \left[
\sum_{2(j+l)<n} \Lambda^{n-2j-2l} b_{2j}(x,x)
\frac{(-m^2)^l l!}{n-2j-2l} \right. \nonumber\\
&&\qquad\qquad\qquad \left.+ \sum_{2(j+l)=n} \ln (\Lambda) 
(-m^2)^l l!\, b_{2j}(x,x) +\mathcal{O}(\Lambda^0) \right]\,.
\label{divWLam}
\end{eqnarray}
We see that the ultra violet divergences in the one-loop effective
action are defined by the heat kernel coefficients $b_{k}(x,x)$
with $k\le n$.

On non-compact manifolds the integral of $b_0(x,x)$ is divergent.
This divergence is removed by subtracting a ``reference'' heat
kernel (see sec.\ \ref{s6sinpot}).

Contributions from higher heat kernel coefficients $b_k$, $k>n$
to the effective action are not divergent and can be easily
calculated yielding in the limit $\Lambda\to\infty$
\begin{equation}
-\frac 12 (4\pi )^{-n/2} m^n \int d^nx \sqrt{g}
\sum_{2j>n} \frac{b_{2j}(x,x)}{m^{2j}} \Gamma (2j-n)\,.
\label{largemass}
\end{equation}
This is nothing else than the large mass expansion of the effective
action. This expansion is valid for relatively weak and slowly varying
background fields.

We have seen that the heat kernel expansion 
describes
\begin{itemize}
\item short-distance behaviour of the propagator;
\item one-loop divergences and counterterms;
\item $1/m$ expansion of the effective action.
\end{itemize}
We shall see below that heat kernel provides a natural framework
for studying 
\begin{itemize}
\item quantum anomalies (sec.\ \ref{sec7});
\item various perturbative expansions of the effective action
(sec. \ref{sec8});
\item selected non-perturbative relations for the effective
action (sec.\ \ref{sec9}).
\end{itemize}

Of course, in all these applications the heat kernel methods have
to compete with other techniques. The main advantage of the heat
kernel is that it delivers necessary information in terms of just
few geometric invariants. This method does not make distinctions 
between different spins, gauge groups, etc. Even dependence of the
space-time dimensionality is in most cases trivial. Therefore, on
one hand, just a single calculation serves then in many applications.
On the other hand, calculations in simple particular cases give
valuable information on the general structure of the heat kernel.
This property is especially useful when one deals with complicated
geometries (like in the presence of boundaries or singularities).
During the last decade, many models which lead to such complicated
geometries were very actively studied in theoretical physics.
The Dirichlet branes and the brane world scenario are the most 
popular but not the only examples.

We have to mention also the limitations of the heat kernel
formalism. It works less effectively in the presence of spinorial
background fields, i.e. when there is mixing between bosonic
and fermionic quantum fields. This problem is probably of the
technical nature, so that the corresponding formalism may be
developed some time in the future. A more serious drawback is
that the heat kernel expansion is not applicable beyond the
one-loop approximation. It is not clear whether necessary
generalisations to higher loop could be achieved at all.  

It is not possible to write a review paper on heat kernel which
would be complete in all respects, especially in the bibliography.
A more comprehensive treatment of many mathematical problems related
to the heat kernel expansion can be found in 
\cite{Gilkey:1995,Kirsten:2001wz}. The book by Kirsten \cite{Kirsten:2001wz}
considers also specific physical applications as the Casimir energy and
the Bose--Einstein condensation. The recent review paper 
\cite{Bordag:2001qi} is devoted to the Casimir effect
(see also \cite{Milton:2001yy}). The monographs
by Birrell and Davies \cite{Birrell:1982ix} and Fulling 
\cite{Fulling:1989nb} remain standard sources on quantum field theory
in curved space. Quantization of gauge theories is explained in
\cite{Faddeev:1980be}. 
In  \cite{Buchbinder:1992rb,Esposito:1997wt,Esposito:1997mw,
Avramidi:2000bm,Esposito:2001rx}
the heat kernel expansion is treated from the point of view of quantum
gravity and quantum cosmology. Useful information about properties of
the zeta function can be found in \cite{Elizalde:1995zk}.
One may also consult \cite{EORBZ}. The DeWitt approach to the heat kernel
and its' generalisations are described in \cite{Barvinsky:1985an}.
The path integral point of view on the heat kernel can be found
in \cite{Roepstorff:1994ga}.

Our primary goal is local formulae for the heat kernel coefficients.
Therefore, in some cases global aspects will be somewhat neglected.

This report is organised as follows. The next section contains
necessary preliminary information on spectral geometry and
differential geometry. There we define main geometric characteristics
of the manifold and of the boundary. We discuss the zeta function
(which defines the effective action) and the resolvent (which is
a generalisation of the propagator). We relate the asymptotics of these
functions to the heat kernel coefficients. Through this report we
work on Euclidean manifolds. A short remark on the
analytical continuation to the Lorentzian signature is given at the
end of section \ref{sec2}.
In  section \ref{sec3} we consider
the most widely used models of quantum field theory and open bosonic
strings. The one-loop dynamics in each of these models is defined by
a second order differential operator which depends on an effective
connection and on a matrix-valued potential. The connection and
the potential serve as a basis of an invariant description of all
that models in the language of spectral geometry. These quantities
are written down explicitly for each model. We also define suitable
boundary conditions. In section \ref{sec4} we consider the heat kernel
expansion on manifolds without boundary. We introduce a simple and
very powerful method and illustrate it by calculating several leading
terms in the heat kernel expansion. We also briefly discuss some other
methods and non-minimal operators. Section \ref{sec5} is devoted to 
manifolds with boundary. The heat kernel expansion for standard
Dirichlet, Neumann, and mixed boundary conditions is considered in some
detail. We also describe less known oblique and spectral boundary   
conditions (these are the ones which contain tangential derivatives
on the boundary or non-local projectors). We discuss loss of the
so-called strong ellipticity for oblique boundary conditions which
corresponds to the critical value of the electric field in string 
physics. With the example of spectral boundary conditions we
illustrate appearance of non-standard ($\ln t$) asymptotics in the
heat kernel. The results of sections \ref{sec4} and \ref{sec5}
are valid on smooth manifolds with smooth potentials and gauge field.
In section \ref{sec6} we consider the case when either the 
background fields or the manifold itself have singularities.
In particular, conical and domain wall singularities are considered.

Sections \ref{sec7} - \ref{sec9} are devoted to applications.
In section \ref{sec7} we relate conformal and chiral anomalies
to certain heat kernel coefficients and re-derive the anomalies
in several particular models. This is a very spectacular but also 
rather well known application of the heat kernel expansion.
In section \ref{sec8} we go beyond
the power series in the proper time $t$. We consider mainly two
particular cases. The first one is the case when derivatives are more important
than the potentials (so that we sum up leading non-localities).
The second one is the so-called low energy expansion which neglects
derivatives of the background fields starting with certain order
(but treats all powers of background potentials and curvatures 
exactly). In that section we also review some results on the heat
kernel on homogeneous spaces where spectrum of relevant operators
may be found exactly. In section \ref{sec9} we consider two examples
when the heat kernel can be used to obtain exact results for the
effective action. The first one is the famous Polyakov action.
The second example is exact relations between the effective actions
in dual theories. In this section (in contrast to section \ref{sec8})
we don't have to neglect derivatives or powers of
the background fields. Section \ref{sec10} contains concluding
remarks.


